\documentclass[letterpaper, onecolumn, draftclsnofoot, 10pt, compsoc]{IEEEtran}
\usepackage{graphicx}
\usepackage{url}
\usepackage{setspace}

\usepackage{geometry}
\geometry{textheight=9.5in, textwidth=7in}

% 1. Fill in these details
%\def \CapstoneTeamName{		The Cleverly Named Team}
\def \CapstoneTeamNumber{		63}
\def \GroupMemberOne{			Corynna Park}
\def \GroupMemberTwo{			Joshua Lioy}
\def \GroupMemberThree{			Jackson Wells}
\def \CapstoneProjectName{		Molecules in 3D?! And in color!? That I can hold in my hand? No way!!!}
\def \CapstoneSponsorCompany{	Oregon State University College of Science Department of Biochemistry and Biophysics}
\def \CapstoneSponsorPerson{		Dr. Victor Hsu}

% 2. Uncomment the appropriate line below so that the document type works
\def \DocType{		Problem Statement
				%Requirements Document
				%Technology Review
				%Design Document
				%Progress Report
				}
			
\newcommand{\NameSigPair}[1]{\par
\makebox[2.75in][r]{#1} \hfil 	\makebox[3.25in]{\makebox[2.25in]{\hrulefill} \hfill		\makebox[.75in]{\hrulefill}}
\par\vspace{-12pt} \textit{\tiny\noindent
\makebox[2.75in]{} \hfil		\makebox[3.25in]{\makebox[2.25in][r]{Signature} \hfill	\makebox[.75in][r]{Date}}}}
% 3. If the document is not to be signed, uncomment the RENEWcommand below
%\renewcommand{\NameSigPair}[1]{#1}

%%%%%%%%%%%%%%%%%%%%%%%%%%%%%%%%%%%%%%%
\begin{document}
\begin{titlepage}
    \pagenumbering{gobble}
    \begin{singlespace}
    	\includegraphics[height=4cm]{coe_v_spot1}
        \hfill 
        % 4. If you have a logo, use this includegraphics command to put it on the coversheet.
        %\includegraphics[height=4cm]{CompanyLogo}   
        \par\vspace{.2in}
        \centering
        \scshape{
            \huge CS Capstone \DocType \par
            {\large\today}\par
            \vspace{.5in}
            \textbf{\Huge\CapstoneProjectName}\par
            \vfill
            {\large Prepared for}\par
            \Huge \CapstoneSponsorCompany\par
            \vspace{5pt}
            {\Large\NameSigPair{\CapstoneSponsorPerson}\par}
            {\large Prepared by }\par
            Group\CapstoneTeamNumber\par
            % 5. comment out the line below this one if you do not wish to name your team
            %\CapstoneTeamName\par 
            \vspace{5pt}
            {\Large
                \NameSigPair{\GroupMemberOne}\par
                \NameSigPair{\GroupMemberTwo}\par
                \NameSigPair{\GroupMemberThree}\par
            }
            \vspace{20pt}
        }
        \begin{abstract}
        % 6. Fill in your abstract    
        	Our project aims to produce a robust workflow for printing 3D multi-color models. Currently, the capability to produce monochromatic 3D models already exists. We look to expand upon this by adding a multi-color functionality. This will enable enhanced learning opportunities and better absorption of concepts, specifically those relating to the field of biochemistry. While this project originated as one for the OSU Biochemistry department, we hope that our findings about how to perform the process of multi-color 3D printing will benefit those looking to print 3D multi-color models in general, not just those relating to Biochemistry.\par
            To complete this project, we will need to become familiar with the existing workflow of how to print monochromatic models. From there, we will utilize existing free and open-source software to deduce the best way to create multi-color 3D models. At the completion of this project, we will have produced a robust workflow as well as a journal article describing the process we created to utilize 3D multi-color printing.\par
        \end{abstract}     
    \end{singlespace}
\end{titlepage}
\newpage
\pagenumbering{arabic}
\tableofcontents
% 7. uncomment this (if applicable). Consider adding a page break.
%\listoffigures
%\listoftables
\clearpage

% 8. now you write!
\section{Project Definition}
    \begin{singlespace}
        Tactile learning is an effective tool for instructors help improve student understanding of concepts and ideas. Biochemistry courses involve countless unique structures all with unique functions so 3D printing allows teachers to take learning one step further. The Biochemistry department at Oregon State University recently purchased two 3D printers, the Zmorph 2.0S and the Rostock MAX delta, that are able to print models in larger numbers at a reasonable cost. While the current setup is adequate, the printers are unable to print models with multiple colors, creating objects that don’t convey the information that they need to.\par
        The department has now procured a multifilament printing interface, the Palette+ by Mosaic Manufacturing, which opens the door to multicolor 3D models. The Palette+ works with both printers and is able to produce prints of up to 4 colors natively, and prints of 8 colors are claimed to be possible \cite{8colorprint}. This device greatly expands the realm of potential models to be produced by the department but the technology is so new, the creation of 3D multicolor files is a tedious process. No clear pipeline for file creation currently exists and many of the software tools involved lack proper documentation or are extremely expensive. Multiple publications also exist on the topic of 3D printing but multicolor prints go unmentioned \cite{progconv}.\par
        The department is seeking a robust and streamlined workflow to use this technology to its full potential. While mostly for biochemical structure prints, the workflow should be able to handle any structure desired.\par
    \end{singlespace}

\section{Proposed Project Solution}
    \begin{singlespace}
        Our proposed solution is to develop a robust workflow for multicolor 3D printing. Additionally, to maximize the benefit of our project accomplishments to the public, provide online documentation of the multi-color 3D printing process by writing a journal article describing the process and then submitting it for publication. After we become well informed about the monochromatic 3D printing process, we will need to explore the various options we have regarding software tools to generate printable 3D object files. To address the problem at a software level there are several tools that could potentially be used to implement the necessary workflow. Firstly, making use of existing open source or freely available tools to input usable color data into the 3D object files. These files then need to be “sliced” so that the object data can be interpreted by the 3D printer. This is done by converting the object files into gcode, which is a language that the 3D printer uses to print a given object.\par
        At this point in the process, the gcode needs to be processed so that the Palette+ can determine which filament to feed to the 3D printer at any given time. This is where it may be appropriate for us to create a proprietary tool that interprets the gcode and sends the correct commands to the Palette+. This could be done through scripting or other means depending on what fits best into the workflow that we design and implement. We will need to decide the best file type to use to show and produce a multi-colored model, as multiple file types exist.\par
    \end{singlespace}
    
\section{Performance Metrics}
    \begin{singlespace}
        At the completion of this project the program created will be able to be operated without confusion or informational aid.
        The program will be able to produce prints consisting of at least 4 fours colors, given any object.
        Colors on the object should not mix, and all lines should be well defined.
        Biological structure files (.pdb, .cif, etc.) can be taken as input, processed/colored, and printed on the supplied 3D printers.
        Non-biological files will be supported, including all the standard 3D object file types, to allow for printing of any model desired by the user. A GitHub page will be available to the public to accompany the publication.\par
    \end{singlespace}

\section{Conclusion}
    \begin{singlespace}
        In conclusion, our project stands to benefit several parties, from the educator to the curious. We look forward to being able to help kinesthetic learners learn to their best ability thanks to the advancements in 3D printing.\par
        With the ability to 3D print models in color, individuals will be able to have a better physical grasp on many concepts, and not just concepts limited to the field of biochemistry. We are not just creating a workflow for chemists and biologists, or instructors and students, we are creating a workflow to create a greater desire for 3D printing and the benefits that they can bring.\par
        This project aims to improve an already impressive and useful technology and make it easier to use by providing a streamlined and proven process. As such, it will simultaneously improve user experience and achievement regarding 3D printing.\par
    \end{singlespace}

\bibliography{references}{}
\bibliographystyle{ieeetr}

\end{document}
